\section{Spherical probabilities and exact observations}
\label{p:geometry-model}

We collect the geometric facts used by all three estimates and then specify
the finite-state experiment to which they will be applied. The geometric
lemmas keep their normalizations visible: all measures on spheres below are
probability measures, whereas densities on Euclidean label spaces are
relative to Lebesgue measure unless another reference measure is stated.

\subsection{Coordinates and caps}

Write \(\sigma_d\) for uniform probability on the unit sphere
\(S^{d-1}\subset\R^d\), and write \(B(z,r)\) for the closed Euclidean ball
with center \(z\) and radius \(r\). When the ambient dimension is fixed we
abbreviate \(\sigma_d\) to \(\sigma\). The norm is Euclidean, and
\(\norm A_{\op}\) denotes the operator norm of a matrix. Constants denoted
by \(C\) are positive and absolute and may increase from one occurrence to
the next. A subscript records any permitted additional dependence.
All logarithms are natural unless a base is displayed.

\begin{lemma}[Coordinates of a uniform direction]\label{p:coordinate-density}
Let \(D\ge2\) and \(1\le k<D\) be integers. If \(U\) is uniform on
\(S^{D-1}\) and \(P\) is any \(k\)-by-\(D\) matrix satisfying
\(PP^{\mathsf T}=I_k\), then \(PU\) has Lebesgue density
\begin{equation}\label{p:coordinate-formula}
 c(D,k)(1-\norm u^2)^{(D-k-2)/2}\1_{\{\norm u<1\}},
 \qquad
 c(D,k)=\frac{\Gamma(D/2)}
 {\pi^{k/2}\Gamma((D-k)/2)}.
\end{equation}
The direction of the component of \(U\) in \(\ker P\) is uniform on its
unit sphere and is independent of the pair consisting of \(PU\) and the
length of that complementary component. In particular, its conditional
law given \(PU\) is uniform.
\end{lemma}

\begin{proof}
Rotation invariance reduces the projection to the first \(k\) coordinates.
Let \(G\in\R^k\) and \(H\in\R^{D-k}\) be independent standard Gaussian
vectors. Then \((G,H)/\sqrt{\norm G^2+\norm H^2}\) is uniform on the sphere.
The random variables \(\norm G^2/2\) and \(\norm H^2/2\) are independent
unit-scale gamma variables of shapes \(k/2\) and \((D-k)/2\). Changing
variables from these two positive numbers to their sum and their first
fraction shows that the fraction has density
\[
 \frac{\Gamma(D/2)}
 {\Gamma(k/2)\Gamma((D-k)/2)}
 t^{k/2-1}(1-t)^{(D-k)/2-1},\qquad 0<t<1.
\]
The direction of \(G\) is uniform and independent of these lengths.
Polar coordinates in \(\R^k\), with
\(\vol(S^{k-1})=2\pi^{k/2}/\Gamma(k/2)\), now give
\eqref{p:coordinate-formula}. The surface-area identity itself follows by
integrating \(e^{-\norm x^2}\) in polar coordinates. The direction of \(H\)
is uniform and independent of both lengths and of \(G\), which proves the
conditional assertion. The event \(H=0\) is null.
\end{proof}

The beta distribution in this lemma is a classical description of random
orthogonal projections; see Frankl and Maehara~\cite{FranklMaehara1990}.
We use the displayed derivation, including its conditional direction law.

Write \(v_k\) for the Lebesgue volume of the unit ball in \(\R^k\).
The following estimate will cancel the dimensional density scale in both
projection arguments.

\begin{lemma}[Unit-ball volume]\label{p:unit-ball-volume-lemma}
For every integer \(k\ge1\),
\begin{equation}\label{p:unit-ball-volume}
 v_k\le(2\pi e/k)^{k/2}.
\end{equation}
\end{lemma}

\begin{proof}
On the ball of radius \(\sqrt k\), the function
\(e^{-\norm z^2/2}\) is at least \(e^{-k/2}\). Comparing its integral
there with the full Gaussian integral gives
\[
 e^{-k/2}k^{k/2}v_k
 \le\int_{\R^k}e^{-\norm z^2/2}\,dz=(2\pi)^{k/2},
\]
which is the claimed bound.
\end{proof}

\begin{lemma}[Spherical cap mass]\label{p:cap-mass}
Put \(n=d-1\) for \(d\ge2\). There is an absolute \(a>0\) such that,
for every \(c\in S^{d-1}\) and \(0<r\le2\),
\begin{equation}\label{p:cap-mass-two-sided}
 (ar)^n\le \sigma\bigl(B(c,r)\bigr)\le\min\{1,r^n\}.
\end{equation}
For every ambient center \(z\in\R^d\) and radius \(r>0\),
\begin{equation}\label{p:ambient-cap-mass}
 \sigma\bigl(B(z,r)\bigr)\le\min\{1,(2r)^n\}.
\end{equation}
An output direction independent of the signal has probability at most
\(\epsilon^n\) of angular error at most \(\epsilon\), averaged over
\(\sigma\).
\end{lemma}

\begin{proof}
For \(r\le1\), rotate \(c\) to the north pole. A sphere point in
\(B(c,r)\) has last coordinate at least \(1-r^2/2\ge1/2\). Its surface
element as a graph over \(c^\perp\) is at most twice Lebesgue measure,
and its projection lies in the \(n\)-ball of radius \(r\). The full sphere
has area at least twice the volume of the unit \(n\)-ball, by projecting
both hemispheres. This proves the upper bound for \(r\le1\); total mass
one proves it for \(r>1\).

For the lower bound when \(r\le1\), the angular radius
\(2\arcsin(r/2)\) is at least \(r\). In polar angle the cap mass is the
integral of \(\sin^{n-1}\theta\) over this interval divided by its
integral over \([0,\pi]\). Use
\(\sin\theta\ge2\theta/\pi\) for \(0\le\theta\le r\) and bound the
denominator by \(\pi\). Since \(n\le2^{n-1}\), the result is at least
\((r/\pi)^n\). For \(1\le r\le2\), monotonicity therefore gives the
lower bound with \(a=1/(2\pi)\).

If an ambient ball meets the sphere, choose \(c\) in the intersection.
Its sphere portion is contained in \(B(c,2r)\), proving
\eqref{p:ambient-cap-mass}; an empty intersection is immediate. Finally,
angular error at most \(\epsilon\) implies chordal error at most
\(\epsilon\), because \(2\sin(\theta/2)\le\theta\). Apply the
sphere-centered upper bound to each fixed output and then integrate over
any independent output randomization.
\end{proof}

\subsection{The finite-state regression experiment}

An unknown direction \(s\in S^{d-1}\) is observed through independent
standard Gaussian features and their exact labels:
\[
 X_t\sim N(0,I_d),\qquad Y_t=\inner{X_t}{s},\qquad t=1,2,\ldots .
\]
At each sample index the learner retains a state from a set of at most
\(2^M\) values, where \(M\) is a nonnegative integer. A transition may
use the current state, the entire current pair \((X_t,Y_t)\), the sample
index, and fresh independent randomness. Computation within the transition
is unrestricted. Every piece of information depending on earlier samples
that remains available later must be encoded by the finite state. The
learner neither chooses a feature nor revisits an earlier one.

The learner stops after at most a prescribed integer \(T\ge0\) samples.
Its stopping decisions obey the same information restriction. For fixed
rules, its output is a unit vector whose law depends only on the terminal
state and stopping index, with fresh randomness allowed. In particular,
the last observed pair influences the output only through the terminal
state. The sample index is available without being stored in the state;
the proofs below account for the information in the stopping index.

A shared random seed may select the rules before sampling. The pair
consisting of this seed and the initial finite state is jointly independent
of the signal and the entire sequence of feature rows. The seed and initial
state may depend on each other. Fresh transition and output randomness is
independent of these variables and of initialization. Nonuniform rules may
also depend on \(d,\epsilon\), and the sample index. This is a bit-memory
model with a deterministic finite horizon.

We impose the following measurability convention. The original experiment
is jointly measurable in the signal, pre-generated feature rows, shared
seed, initial state, and fresh randomness, including the resulting state
path, stopping index, and output. For almost every fixed shared seed, the
finite vector of transition and stopping probabilities at each fixed
state and index is measurable in the completion of
\(\gamma_d(dx)\,dy\), where \(\gamma_d=N(0,I_d)\). A rule given instead
as a completed-measurable function of the current pair and an independent
fresh seed is also allowed; Lemma~\ref{p:borel-kernels} justifies this
presentation. Each terminal output law is a probability measure on the
Borel sphere.

For the lower bounds, let \(S\sim\sigma\) be independent of the feature
rows. We call the probability of
\(\arccos\inner{\widehat S}{S}\le\epsilon\), including all learner
randomness, the \emph{uniform-sphere average success probability}.

\begin{theorem}[Finite-memory precision bound]\label{p:regression}
Let \(M(d)=o(d^2)\) be a sequence of nonnegative integers and let
\(0<\epsilon(d)\le1/10\). There is an absolute \(c>0\) such that, for
all sufficiently large \(d\), every learner in the experiment above whose
uniform-sphere average success satisfies
\[
 \Prob\{\arccos\inner{\widehat S}{S}\le\epsilon(d)\}\ge\frac23
\]
must use
\[
 T(d)\ge c\,d\log\frac1{\epsilon(d)}.
\]
The dimension threshold may depend on the memory sequence. In particular,
the conclusion holds when the same success probability is required
separately for every \(s\in S^{d-1}\).
\end{theorem}

An every-signal guarantee implies the average guarantee by integration
against \(\sigma\). All three proofs use the average guarantee itself.
The first proof also yields a uniform dependence on the memory budget.

\begin{corollary}[Finite-dimensional memory dependence]
\label{p:finite-memory-tradeoff}
There are absolute constants \(c>0\) and \(d_0\) such that, for every
integer \(d\ge d_0\), every nonnegative integer \(M\), and every
\(0<\epsilon\le1/10\), a learner in the experiment above with
uniform-sphere average success at least \(2/3\) and deterministic finite
horizon \(T\) satisfies
\[
 T\ge \frac{c\,d\log(1/\epsilon)}{1+M/d^2}.
\]
\end{corollary}

The proof appears at the end of Section~\ref{p:projection}. In particular,
the same precision bound holds with absolute constants throughout
\(M\le d^2\) once \(d\) exceeds an absolute threshold.

\subsection{Measurable continuations and shared randomness}

The block arguments will start fresh executions from specified states,
including states the learner never reaches. The next lemma supplies Borel
rules for all such continuations while preserving the actual uniform-prior
experiment. Its proof uses the marginal law of each observation pair,
so it does not require a density conditional on a selected state.

\begin{lemma}[Borel versions and counterfactual continuations]
\label{p:borel-kernels}
Fix \(d\ge2\), finite \(T\), and finite state sets. Suppose each transition
and stopping probability, for each fixed state and index, is measurable in
the completion of the Borel sigma field of \(\R^d\times\R\) under
\(\lambda(dx,dy)=\gamma_d(dx)\,dy\), where \(\gamma_d\) is standard
Gaussian probability on \(\R^d\). The terminal output law depends only on
the terminal state and index. There are everywhere Borel probability
kernels with the same state sets and output laws that induce the same
uniform-prior joint law of signal, state path, stopping index, and output.
For the Borel learner, starting from any state and index and supplying fresh
future features defines a measurable continuation success function, whether
or not that state is reachable in the original execution.
The same conclusion holds when a rule is given before seed integration by
a finite probability vector measurable in the completion of
\(\lambda\otimes\kappa\), where \(\kappa\) is the law of an independent
fresh seed.
\end{lemma}

\begin{proof}
First suppose the rule includes a fresh seed of law \(\kappa\). Each
completed-product measurable coordinate has a product-measurable version:
approximate it by simple functions and replace the countably many level
sets by product-measurable sets modulo null sets. Replace the resulting
finite vector by a fixed point mass wherever it fails to be a probability
vector. This repair occurs on a product-measurable null set. Integrating
against \(\kappa\) now gives a Borel probability vector in the current
pair, and Fubini shows agreement with the original integrated rule for
\(\lambda\)-almost every pair. It therefore suffices to consider the
coordinate probabilities in the first statement.

For each fixed state and index, combine the possible transitions and stops
into a finite set of destinations, tagged as continuing or stopping.
Choose Borel versions of their coordinate probabilities by the same
simple-function construction. On the Borel null set where the versions
fail to be nonnegative or to sum to one, replace the vector by a fixed
point mass. Finitely many states and indices are involved, so all changes
lie in one Borel \(\lambda\)-null set \(Z\).

For every nonzero feature \(x\), Lemma~\ref{p:coordinate-density} gives
a Lebesgue density for \(\inner{x}{S}\). Since a Gaussian feature is
nonzero almost surely, every pre-generated pair
\((X_t,\inner{X_t}{S})\) has law absolutely continuous with respect to
\(\lambda\). Thus each pair avoids \(Z\) almost surely. Even if the
current state depends on previous observations, the probability of being
in a particular state and seeing a pair in \(Z\) is bounded by the
marginal probability of seeing that pair in \(Z\), which is zero.

Almost surely none of the finitely many pre-generated pairs lies in
\(Z\). Couple the original and modified finite destination kernels by
the same fresh uniform random numbers. Induction gives identical state
paths and stopping decisions; couple their unchanged terminal output laws
identically as well. This proves equality of the specified joint laws.
Finally, finite composition and integration of the Borel kernels against
independent future features and coins defines a measurable continuation
from every specified state and index, including unreachable states.
\end{proof}

\paragraph{Convention for the three applications.}
We first condition on a shared seed for which the model assumptions hold.
The joint independence of seed and initial state from signal and feature
rows ensures that the conditional initial-state law still contains no
signal or sample information. We then use the Borel learner supplied by
Lemma~\ref{p:borel-kernels}. The lemma need only preserve the
uniform-prior law; preservation at every fixed signal is unnecessary.
All bounds below are uniform in the fixed rules and are averaged over the
shared seed at the end. The original conditional success probabilities
are measurable in that seed because the original experiment is jointly
measurable. Neither the replacement kernels nor any subsequent proof
decompositions need to be selected jointly across seeds.

\subsection{A common full-data endpoint}

The three moment arguments produce bounds involving \(T+d\). The next
elementary obstruction supplies the linear lower bound on \(T\) needed
to remove that additive \(d\), without using any memory restriction.

\begin{lemma}[Antipodal full-data bound]\label{p:haar-linear}
In the experiment above, if \(\epsilon\le1/10\) and \(T\le d/8\),
even an estimator given all \(T\) features and labels has uniform-prior
success probability at most \(5/8\).
\end{lemma}

\begin{proof}
Let \(F\) be the row span of the pre-generated feature matrix and put
\(u=P_FS\). Its dimension is \(T\) almost surely, with the evident
interpretation for \(T=0\). The labels determine \(u\).
For \(T=0\), the residual is the original uniform signal. For \(T>0\),
Lemma~\ref{p:coordinate-density}, applied after fixing the feature matrix,
shows that conditional on all the data the residual is
\(\sqrt{1-\norm u^2}\,\omega\), with \(\omega\) uniform on the unit
sphere of \(F^\perp\). Before conditioning on the labels, isotropy gives
\(\E\norm u^2=T/d\). Consequently
\(\Prob\{\norm u^2\ge1/2\}\le2T/d\).

When \(\norm u^2<1/2\), opposite residual directions give two signals
at distance \(2\sqrt{1-\norm u^2}>\sqrt2>2\epsilon\). They have the same
conditional law, and at most one can be within chordal distance
\(\epsilon\) of a fixed output. Angular success implies this chordal
condition, so conditional success is at most \(1/2\), also after
averaging independent estimator randomness. Total success is at most
\[
 \frac12+\frac12\Prob\{\norm u^2\ge1/2\}
 \le\frac12+\frac Td\le\frac58.
\]
Giving the estimator all pre-generated data also covers a learner that
stops at a prefix, with its independent shared and fresh randomness.
\end{proof}
